% The multiplicity properties of carbon-rich Wolf-Rayet stars in the LMC
% Guy Shtainer & Tomer Shenar — Tel Aviv University
% A&A journal format

\documentclass{aa}

% ============================================================
% Packages (only those NOT already loaded by aa.cls)
% ============================================================
% inputenc, amsmath, amssymb, natbib, xspace already loaded by aa.cls
\usepackage{graphicx}
\usepackage{booktabs}
\usepackage{multirow}
\usepackage{hyperref}

% ============================================================
% Key number macros — single source of truth
% ============================================================
% Observation programme
\newcommand{\esoProgrammes}{105.205Y, 106.213X, 107.22A1\xspace}  % ESO programme IDs (P105--P107)
\newcommand{\Ttotalhrs}{45\xspace}                          % total awarded VLT time [hours]
\newcommand{\Nbright}{20\xspace}                            % bright (V<15) targets
\newcommand{\Nfaint}{5\xspace}                              % faint (V>15) targets
\newcommand{\snrTarget}{70\xspace}                          % S/N target at 450 nm
\newcommand{\Rxshooter}{5\,000--10\,000\xspace}             % X-SHOOTER resolving power
\newcommand{\slitUVB}{0.8\arcsec\xspace}
\newcommand{\slitVIS}{0.9\arcsec\xspace}
\newcommand{\slitNIR}{0.9\arcsec\xspace}

% Sample
\newcommand{\Nstars}{25\xspace}          % apparently single stars analyzed
\newcommand{\Nknown}{3\xspace}           % known binaries from Bartzakos (2001)
\newcommand{\Ntotal}{28\xspace}          % total sample
\newcommand{\Nepochs}{6\xspace}          % typical epochs per star

% Binary detection
\newcommand{\Ndetected}{10\xspace}       % newly detected binaries
\newcommand{\Nbin}{13\xspace}            % total binaries (detected + known)
\newcommand{\fbinobs}{46\%\xspace}       % observed binary fraction (text-mode only -- never wrap in $...$)
\newcommand{\drvthresh}{45.5\xspace}     % delta-RV threshold [km/s]
\newcommand{\sigthresh}{4\xspace}        % significance threshold [sigma]
% Use aa.cls's \ion{}{} macro and \ensuremath for the lambdas so the macro is
% safe to use in text mode, math mode, and section/paragraph titles.
\newcommand{\primaryline}{\ion{C}{iv}\,\ensuremath{\lambda\lambda 5808\text{--}5812}\xspace}

% Bias correction (placeholders — update when errors computed)
\newcommand{\fbincorr}{TBD\xspace}       % corrected binary fraction
\newcommand{\fbinerrup}{TBD\xspace}      % +error
\newcommand{\fbinerrdn}{TBD\xspace}      % -error
\newcommand{\pibestfit}{TBD\xspace}      % best-fit period exponent
\newcommand{\pierrup}{TBD\xspace}        % +error
\newcommand{\pierrdn}{TBD\xspace}        % -error

% Simulation
\newcommand{\Nsim}{10\,000\xspace}       % Monte-Carlo draws per grid point
\newcommand{\sigmasingle}{15\xspace}     % intrinsic RV scatter [km/s]
\newcommand{\sigmameas}{5\xspace}        % measurement uncertainty [km/s]

% Threshold-curve fit (placeholders — update when fit is run)
\newcommand{\fbinThresh}{TBD\xspace}        % best-fit binary fraction from threshold-curve method
\newcommand{\fbinThreshErrUp}{TBD\xspace}   % +error
\newcommand{\fbinThreshErrDn}{TBD\xspace}   % -error
\newcommand{\sigmaIntrinsic}{TBD\xspace}    % fitted intrinsic single-star scatter [km/s]
\newcommand{\sigmaMeasFit}{TBD\xspace}      % fitted per-epoch measurement uncertainty [km/s]

% Spectroscopic morphology investigation (placeholders)
\newcommand{\NbinSpec}{TBD\xspace}          % additional binaries flagged by morphology
\newcommand{\NbinFinal}{TBD\xspace}         % \Nbin + \NbinSpec, total binaries after spectrum pass
\newcommand{\fbinFinal}{TBD\xspace}         % final observed binary fraction (post-spectrum)
\newcommand{\NcandidateBH}{2\xspace}        % WC + compact-object follow-up candidates

% CCF + line-analysis macros (some are derived; some are TBD until final settings locked)
\newcommand{\NlinesAll}{11\xspace}                  % all emission lines initially considered
\newcommand{\fitFracDefault}{0.97\xspace}           % default CCF fit fraction
\newcommand{\Nstarsmincorr}{8\xspace}               % min sample-overlap for line-agreement Pearson
\newcommand{\sigmaSingleFit}{TBD\xspace}            % fitted single-star sigma from parametric Gaussian threshold model
\newcommand{\sigmaBinaryFit}{TBD\xspace}            % fitted binary sigma from same model
\newcommand{\fbinAnalytic}{TBD\xspace}              % f_bin from parametric Gaussian fit (not the MC-corrected fbincorr)

% Physical
% Using \ensuremath so these work correctly in both text and math mode
% without nested $...$ toggles that can cause "small + glued" rendering bugs.
\newcommand{\kms}{\ensuremath{\mathrm{km\,s^{-1}}}\xspace}
\newcommand{\Msun}{\ensuremath{M_\odot}\xspace}
\newcommand{\Lsun}{\ensuremath{L_\odot}\xspace}
\newcommand{\Zsun}{\ensuremath{Z_\odot}\xspace}
\newcommand{\Mdot}{\ensuremath{\dot{M}}\xspace}

% ============================================================
% Title and authors
% ============================================================
\title{The multiplicity properties of carbon-rich Wolf-Rayet stars
       in the Large Magellanic Cloud}

\author{G.\,Shtainer\inst{1}
  \and T.\,Shenar\inst{1}
}

\institute{School of Physics and Astronomy, Tel Aviv University,
           Tel Aviv 6997801, Israel\\
           \email{guyshtainer@mail.tau.ac.il}
}

\date{Received ???; accepted ???}

% ============================================================
% Abstract — A&A structured format: 5 positional arguments
% \abstract{Context}{Aims}{Methods}{Results}{Conclusions}
% ============================================================
\abstract
{Wolf-Rayet (WR) stars of the carbon sequence (WC/WO) represent the final
stripped stages of massive stellar evolution, yet the role of binary
interactions in their formation remains debated. The Large Magellanic Cloud
(LMC), at $Z \approx 0.3$--$0.5\,\Zsun$, provides a critical low-metallicity
environment to test theoretical predictions of the WR binary fraction.}
%
{We aim to determine the spectroscopic binary fraction of LMC WC/WO stars
and constrain the intrinsic binary fraction and orbital period distribution
through Monte-Carlo bias correction.}
%
{We analyse multi-epoch VLT/X-SHOOTER spectroscopy of \Nstars apparently
single WC/WO stars in the LMC across \Nepochs epochs. Radial velocities are
measured via cross-correlation of the \primaryline emission line. Binary
classification requires $\Delta\mathrm{RV} > \drvthresh$\,\kms and a
$\sigthresh\sigma$ significance threshold. The intrinsic binary fraction and
period-distribution exponent are constrained by comparing the observed
$\Delta\mathrm{RV}$ distribution to Monte-Carlo simulations via a
multinomial likelihood comparison (Dsilva et al. 2023).}
%
{We detect \Ndetected new spectroscopic binaries, yielding a total of
\Nbin/\Ntotal (\fbinobs) when combined with \Nknown previously known
binaries. The bias-corrected intrinsic binary fraction is
$f_\mathrm{bin} = {\fbincorr}^{+\fbinerrup}_{-\fbinerrdn}$, with a
period-distribution exponent $\pi = {\pibestfit}^{+\pierrup}_{-\pierrdn}$.}
%
{} % Conclusions — to be finalised once bias-correction errors are settled

\keywords{stars: Wolf-Rayet -- binaries: spectroscopic -- stars: massive -- Magellanic Clouds}

% ============================================================
\begin{document}

\maketitle

% ============================================================
% Main body
% ============================================================
\section{Introduction}
\label{sec:intro}

% --- WR stars as endpoints of massive evolution ---
Wolf-Rayet (WR) stars are evolved, hydrogen-depleted massive stars
characterised by broad emission lines arising from their powerful,
optically thick stellar winds \citep{Crowther2007}. Classical WR stars
(cWR) have shed their hydrogen envelopes, exposing helium-burning cores
with typical masses of $8$--$30$\,\Msun and mass-loss rates of
$\Mdot \sim 10^{-5}$--$10^{-4}$\,\Msun\,yr$^{-1}$. The WR phase
represents one of the final stages before core collapse, making WR stars
direct progenitors of Type~Ibc supernovae and potential progenitors of
stellar-mass black holes \citep{Sander2019}.

% --- WC/WO classification ---
The carbon-sequence WR stars (WC and WO subtypes) are distinguished by
strong carbon and oxygen emission lines, products of core helium burning
brought to the surface by mass loss. They are more chemically evolved than
the nitrogen-sequence (WN) stars, which display CNO-cycle products from
hydrogen burning. The WC/WO phase is therefore a later evolutionary stage,
and these stars are closer to core collapse.

% --- Binary formation channel ---
Two primary channels can produce WR stars: (1)~intrinsic mass loss driven
by radiation pressure on spectral lines, stripping the hydrogen envelope of
sufficiently massive single stars; and (2)~binary interactions, including
Roche-lobe overflow and common-envelope evolution, which strip the envelope
at lower initial masses \citep{Shenar2020}. The relative contribution of
each channel is a fundamental open question in massive star evolution.

% --- Metallicity dependence ---
At solar metallicity, line-driven winds can produce WR stars from single
stars with initial masses $M_\mathrm{ini} \gtrsim 20$\,\Msun. At lower
metallicity, weaker winds raise this threshold ($\gtrsim 25$\,\Msun in the
LMC, $\gtrsim 37$\,\Msun in the SMC), implying that binary interactions
should become increasingly important for WR production
\citep{Shenar2020}. Population synthesis models predict binary fractions
$\gtrsim 50\%$ for WR stars at sub-solar metallicity
\citep{Marchant2023}, but observational constraints remain sparse.

% --- Previous surveys ---
The binary fraction of WR stars has been measured primarily in the Milky
Way. \citet{Dsilva2020} and \citet{Dsilva2023} conducted spectroscopic
multiplicity surveys of Galactic WC/WO and WN stars using
cross-correlation techniques, finding observed binary fractions of
$\sim 40\%$. In the Magellanic Clouds, \citet{Bartzakos2001} surveyed
\Ntotal WC/WO stars in the LMC, detecting only \Nknown spectroscopic
binaries. \citet{Shenar2019} studied the WN stars in the LMC. A
comprehensive survey of the LMC WC/WO population with modern
instrumentation has been lacking.

% --- This work ---
In this paper, we present a spectroscopic multiplicity survey of \Nstars
apparently single WC/WO stars in the LMC, using multi-epoch VLT/X-SHOOTER
observations. We apply the cross-correlation technique of
\citet{Zucker1994} to measure radial velocities and classify binary
candidates. We then perform a Monte-Carlo bias correction following the
methodology of \citet{Dsilva2023} to constrain the intrinsic binary
fraction and orbital period distribution, providing the first such
measurement at LMC metallicity for carbon-sequence WR stars.

The paper is structured as follows. Section~\ref{sec:obs} describes the
sample and observations. Section~\ref{sec:methods} presents the
radial-velocity measurement and binary-classification methodology.
Section~\ref{sec:results} reports the detected binaries and observed binary
fraction. Section~\ref{sec:bias} describes the bias-correction simulation.
Section~\ref{sec:discussion} discusses the results in the context of WR
formation channels and metallicity dependence.
Section~\ref{sec:conclusions} summarises our findings.

\section{Sample and observations}
\label{sec:obs}

We present multi-epoch optical and near-infrared spectroscopy of \Ntotal
carbon-sequence Wolf-Rayet (WC/WO) stars in the Large Magellanic Cloud,
obtained as part of ESO observation programme \esoProgrammes
(PI: T.\,Shenar) under VLT cycles P105--P107. Each target was observed
in \Nepochs epochs, distributed in service mode across three consecutive
semesters, for a total of \Ttotalhrs\,h of awarded VLT time. The cadence
and signal-to-noise targets were tailored to the dual goal of detecting
spectroscopic binaries among nominally single WC stars and of constraining
their underlying orbital period distribution.

\subsection{Target sample}
\label{sec:sample}

Our sample consists of \Ntotal WC- and WO-type Wolf-Rayet stars in the
Large Magellanic Cloud, drawn from the catalogue of
\citet{Breysacher1999} and the survey of \citet{Bartzakos2001}. Of these,
\Nknown were previously identified as spectroscopic binaries by
\citet{Bartzakos2001}; the remaining \Nstars stars were classified as
apparently single and form the primary targets of our survey.
Table~\ref{tab:sample} lists the target stars with their designations,
spectral types, and number of observed epochs. To balance signal-to-noise
requirements against the available VLT time, the sample was stratified
by apparent brightness into \Nbright bright targets ($V < 15$) and
\Nfaint faint targets ($V > 15$), each of which was assigned a distinct
exposure-time and observing-condition strategy (Sect.~\ref{sec:xshooter}).

\begin{figure}
\centering
\includegraphics[width=\columnwidth]{plots/sample_map.pdf}
\caption{Sky positions of the \Ntotal LMC WC/WO targets overlaid on an
  H$\alpha$ image of the Large Magellanic Cloud. The \Nstars apparently
  single targets analysed in this work are shown in blue; the \Nknown
  previously known binaries from \citet{Bartzakos2001} are shown in
  orange.}
\label{fig:sample_map}
\end{figure}

\subsection{VLT/X-SHOOTER spectroscopy}
\label{sec:xshooter}

Multi-epoch spectroscopy was obtained with the X-SHOOTER spectrograph
\citep{Vernet2011} mounted on UT2 (Kueyen) of the ESO Very Large
Telescope at Cerro Paranal, Chile. X-SHOOTER is a three-arm
cross-dispersed echelle providing simultaneous wavelength coverage
across the UVB ($\sim 300$--$560$\,nm), VIS ($\sim 560$--$1020$\,nm),
and NIR ($\sim 1020$--$2480$\,nm) arms. We adopted slit widths of
\slitUVB in the UVB arm, \slitVIS in the VIS arm, and \slitNIR in the
NIR arm, yielding
an intermediate resolving power of $R \sim$ \Rxshooter across the full
wavelength range.

The choice of X-SHOOTER is motivated by two requirements specific to
the spectroscopic identification of WR binaries. First, the broad
emission lines that dominate WC spectra have full widths at half maximum
of order $10^{3}$\,\kms; resolving radial-velocity shifts at the few-\kms
level against this intrinsic width demands a spectral resolution of at
least $R \sim 5\,000$--$10\,000$, as established by previous WR binary
surveys \citep[e.g.][]{Bartzakos2001,Dsilva2022}.
Second, disentangling the WR component from a putative OB-type companion
benefits decisively from a multi-wavelength baseline: characteristic He,
C, and O features are distributed throughout the optical and near-infrared,
and the simultaneous arm-by-arm coverage delivered by X-SHOOTER provides
the broad spectral grasp required to track these features in a single
exposure. X-SHOOTER on the VLT uniquely satisfies both criteria for our
target brightness regime in the LMC.

A further motivation for broad coverage is the well-known sensitivity of
single-line WR radial velocities to wind variability and clumping:
estimates derived from a single emission line are biased by stochastic
line-profile variations of order $10$--$30$\,\kms. Cross-correlating across the eleven distinct He/C/O
emission lines listed in Table~\ref{tab:emission_lines}
(Sect.~\ref{sec:ccf}) averages down this stochastic component and brings
the per-epoch radial-velocity uncertainty into the few-\kms regime,
matching the precision required to detect long-period WR binaries.

The \Nepochs epochs per target were distributed as a 2$+$2$+$2 cadence
across the three observing semesters P105--P107, providing baselines
that range from days to ${\sim}\,500$\,d. This pattern was designed to
retain sensitivity to orbital periods up to $P \sim 10^{3}$\,d for
binaries whose RV semi-amplitudes exceed
$\Delta\mathrm{RV} \approx 30$\,\kms. The signal-to-noise and observing
conditions requested for the two brightness groups differ accordingly.
For the \Nbright bright targets, exposures were configured to reach
$\mathrm{S/N} \geq \snrTarget$ at $450$\,nm under airmass $\leq 2.2$,
seeing $\leq 1.3\arcsec$, and full lunar illumination. For the
\Nfaint faint targets, the requested S/N at $450$\,nm was relaxed to
${\sim}\,50$ but observing conditions were tightened to airmass
$\leq 1.8$, seeing $\leq 1.0\arcsec$, and lunar illumination
$\leq 0.5$, in order to keep the per-epoch read-noise contribution
sub-dominant relative to wind variability.

\subsection{Data reduction}
\label{sec:reduction}

Raw two-dimensional X-SHOOTER frames were reduced with the ESO
X-SHOOTER pipeline, which performs bias subtraction, flat-fielding,
order tracing, sky subtraction, wavelength calibration against ThAr arc
exposures, and flux calibration against spectrophotometric standards,
producing wavelength- and flux-calibrated 2D and 1D spectra per arm.
The subsequent post-processing --- spectral combination, continuum
normalisation, and (where necessary) spatial contamination removal ---
was performed using two custom tools developed for this work, both to
be released as open-source repositories upon publication. The full
post-processing chain is: an initial three-arm stitch
(Sect.~\ref{sec:stitch_initial}); continuum normalisation
(Sect.~\ref{sec:normalisation}); spatial contamination removal applied
arm-by-arm where needed (Sect.~\ref{sec:contamination_removal}); and a
final re-stitch of the cleaned, already-normalised arms into a single
1D per-epoch spectrum (Sect.~\ref{sec:stitch_final}).

\subsubsection{Initial spectral combination}
\label{sec:stitch_initial}

The three spectral arms were combined into a single working spectrum
per epoch using signal-to-noise-squared weighting. In the overlap
regions between adjacent arms, we computed the mean flux in the last
$5\%$ of the overlap for the redder arm and the corresponding region
of the bluer arm, deriving an alignment factor
$A = \bar{F}_\mathrm{red} / \bar{F}_\mathrm{blue}$. The two arms were
considered consistent if
\begin{equation}
|\bar{F}_\mathrm{red} - \bar{F}_\mathrm{blue}| <
\frac{\sqrt{\sigma_\mathrm{red}^2 + \sigma_\mathrm{blue}^2}}{\sqrt{N}}.
\end{equation}
Where the spectral sampling differed (e.g.\ NIR at $0.04$\,nm
vs.\ VIS at $0.02$\,nm), the lower-resolution spectrum was linearly
interpolated onto the finer grid before combination. The combined flux
was computed as $F_\mathrm{comb} = \sum w_i f_i$ with weights
$w_i = \mathrm{SNR}_i^2 / \sum \mathrm{SNR}_i^2$. This initial 1D
combined spectrum per epoch serves as the working spectrum for
continuum-anchor selection (Sect.~\ref{sec:normalisation}).

\subsubsection{Continuum normalisation}
\label{sec:normalisation}

Continuum normalisation was performed with a custom interactive
normalisation tool developed for this work and to be released as an
open-source repository upon publication. Continuum anchor points were
placed by visual inspection in featureless regions of the spectrum,
deliberately avoiding emission and absorption features as well as
telluric residuals. The flux assigned to each anchor was the
iteratively $3\sigma$-clipped mean over a narrow wavelength window
centred on the anchor wavelength, which suppresses the influence of
cosmic rays, hot pixels, and weak unresolved features that would
otherwise displace the local continuum estimate.

The pseudo-continuum was constructed as a piecewise linear interpolation
between consecutive anchor wavelengths, and the observed spectrum was
then divided by this interpolation to yield the normalised flux. The
same anchor wavelengths were reused for all \Nepochs epochs of a given
star, which enforces an epoch-to-epoch consistent normalised continuum
and removes a known source of artificial radial-velocity scatter that
arises when each epoch is normalised independently. Because the anchors
are defined in wavelength space, they can be reapplied unchanged to the
cleaned per-arm spectra produced by the contamination-removal step
(Sect.~\ref{sec:contamination_removal}); this means that no second pass
of anchor placement is required on the re-stitched spectrum
(Sect.~\ref{sec:stitch_final}).

\subsubsection{Spatial contamination removal}
\label{sec:contamination_removal}

In service-mode observations the X-SHOOTER slit can capture, in addition
to the WR target, a serendipitous source projected along the line of
sight that is \emph{not} a binary companion of the WR star --- for
example, an unrelated foreground or background star whose spatial
profile happens to fall within the slit during the exposure. If left
untreated, this contamination biases the 1D-extracted continuum and,
more critically, can imprint absorption features on the WR
emission-line profiles at overlapping wavelengths, mimicking or masking
genuine companion signatures.

To mitigate this, we developed a custom contamination-removal tool that
operates directly on the two-dimensional pipeline products, again to
be released as an open-source repository upon publication. The
identification of the WR's spatial position on the slit is unambiguous
because the WR emission lines are themselves localised tracers of the
WR's row in the spatial direction: their flux is strongly peaked at the
WR's spatial profile centre, and the \primaryline complex in particular
shows up as a bright vertical band in the 2D frame even when the
adjacent contaminant is comparable in continuum brightness.

The user manually selects, for each affected epoch, a spatial-direction
range that brackets the WR profile (typically a few arcseconds in
extent) and excludes the contaminating source. Only flux within this
range is collapsed along the spatial axis to produce the cleaned 1D
spectrum. The procedure is performed independently for each of the
three X-SHOOTER arms, so that the WR's spatial profile is delineated
on the actual 2D frame of each arm rather than inferred from one arm
and propagated to the others; this is important because the spatial
profile, slit illumination, and contaminant brightness ratio all vary
arm-to-arm. The output is one cleaned 1D spectrum per arm per affected
epoch. Stars showing no contamination in any arm are passed through
this step unchanged, with the user-selected spatial range covering the
full slit width, so that the workflow returns the standard pipeline
extraction in those cases.

\subsubsection{Final spectral combination}
\label{sec:stitch_final}

The cleaned per-arm 1D spectra are recombined into a single 1D
spectrum per epoch using the same signal-to-noise-squared weighting
and overlap-region alignment procedure described in
Sect.~\ref{sec:stitch_initial}. Because the continuum anchor wavelengths
defined in Sect.~\ref{sec:normalisation} are identical across all epochs
and are equally valid for the cleaned and uncleaned per-arm spectra,
the re-stitched spectrum is automatically normalised at the same
continuum points; no second round of anchor placement is required on
the final spectrum. The output is a single normalised, contamination-free
1D spectrum per star per epoch.

The resulting set of \Ntotal stars $\times$ \Nepochs epochs of normalised,
cleaned 1D spectra forms the input to the radial-velocity measurement
and binary classification pipeline described in Sect.~\ref{sec:methods}.

\section{Radial-velocity measurement and binary classification}
\label{sec:methods}

\subsection{Cross-correlation function}
\label{sec:ccf}

Radial velocities were measured using the one-dimensional
cross-correlation function (CCF) technique of \citet{Zucker1994}, in the
same form as adopted for the parallel Galactic WC survey of
\citet{Dsilva2020}. For two flux arrays $f_1$ and $f_2$ defined on a
common logarithmic wavelength grid of length $N$, the CCF is the
normalised dot product
\begin{equation}
\rho(s) \;=\; \frac{1}{N\,\sigma_{f_1}\,\sigma_{f_2}}
              \sum_{i=1}^{N} f_1(\lambda_i)\,f_2(\lambda_i + s),
\end{equation}
where $s$ is the velocity shift, and the velocity shift of $f_2$ relative
to $f_1$ is recovered from the position of the CCF maximum. For each star,
a high signal-to-noise template spectrum was constructed by first
computing the CCF of all epochs against an arbitrary reference epoch to
determine relative velocity shifts, correcting each epoch to a common rest
frame, and then co-adding the corrected spectra with signal-to-noise
squared weighting. This iterative procedure removes the dependence on the
choice of reference epoch.

The CCF was computed independently for each of the \NlinesAll emission
lines listed in Table~\ref{tab:emission_lines}, spanning the full
X-SHOOTER wavelength range from the ultraviolet
(\ion{O}{v}\,$\lambda\lambda 3100$--$3175$) to the near-infrared
(\ion{C}{iv}\,$\lambda 20842$). Each line is associated with a default
integration window in wavelength (Table~\ref{tab:emission_lines}) and a
default CCF fit fraction $f_\mathrm{fit} = \fitFracDefault$, which
controls the fraction of the CCF peak used for parabolic fitting: the
fit region extends to the wavelengths at which
$\rho(s) = f_\mathrm{fit}\,\rho_{\max}$, and a parabola of the form
$\rho(s) \simeq a s^2 + b s + c$ is fitted to the points above this
threshold. The radial velocity is then recovered from the parabola apex,
$v = -b/(2a)$, and the formal per-epoch uncertainty
$\sigma_\mathrm{RV}$ is derived from the curvature of the CCF peak
following \citet{Zucker2003},
\begin{equation}
\sigma_\mathrm{RV}
   \;=\; \left[\,
     -\,N \,\frac{\rho''(\hat s)\,\rho(\hat s)}{1 - \rho(\hat s)^2}
   \,\right]^{-1/2},
\end{equation}
where $\hat s$ is the position of the CCF maximum and $\rho''$ is the
second derivative of the CCF at the peak.

For a small fraction of star--line combinations the default settings
yielded a poorly converged parabolic fit -- typically because of edge
effects from the chosen wavelength range, or a noise spike close to the
peak -- and the wavelength range and/or fit fraction were adjusted on a
per-star, per-line basis. All such overrides are documented in the public
CCF settings file released alongside this paper, and their net effect on
the RV measurements is negligible compared to the per-epoch
$\sigma_\mathrm{RV}$.

% Emission line table
\begin{table}
\caption{Emission lines used for cross-correlation analysis.}
\label{tab:emission_lines}
\centering
\begin{tabular}{lcc}
\hline\hline
Line ID & $\lambda_\mathrm{min}$ [nm] & $\lambda_\mathrm{max}$ [nm] \\
\hline
\ion{O}{v}\,$\lambda\lambda 3100$--$3175$   & 310.0 & 317.5 \\
\ion{O}{iv}\,$\lambda\lambda 3350$--$3480$  & 335.0 & 348.0 \\
\ion{C}{iv}\,$\lambda\lambda 3650$--$3900$  & 365.0 & 390.0 \\
\ion{He}{ii}\,$\lambda 4686$                & 456.0 & 480.0 \\
\ion{O}{vi}\,$\lambda\lambda 5210$--$5340$  & 521.0 & 533.5 \\
\ion{He}{ii}\,$\lambda 5412$ \& \ion{C}{iv}\,$\lambda 5471$ & 535.0 & 554.0 \\
\ion{C}{iv}\,$\lambda\lambda 5808$--$5812$  & 570.0 & 588.0 \\
\ion{C}{iii}\,$\lambda\lambda 6700$--$6800$ & 666.5 & 684.0 \\
\ion{C}{iv}\,$\lambda 7063$                 & 697.0 & 714.0 \\
\ion{C}{iv}\,$\lambda 17396$                & 1710.0 & 1763.0 \\
\ion{C}{iv}\,$\lambda 20842$                & 2050.0 & 2100.0 \\
\hline
\end{tabular}
\end{table}

\subsection{Choice of primary RV line}
\label{sec:primary_line}

At the outset of the analysis we made no \emph{a priori} commitment to
which of the \NlinesAll emission lines would yield the most reliable
radial velocity. We therefore measured the CCF for all of them, and
selected the primary RV line based on the physical and empirical
considerations set out below.

\paragraph{NIR lines and telluric H$_2$O contamination.}
The two near-infrared carbon features in our line list,
\ion{C}{iv}\,$\lambda 17396$ and \ion{C}{iv}\,$\lambda 20842$, fall on
the red wing of the strong $2\nu_2 + \nu_3$ / $\nu_1 + \nu_3$ H$_2$O
combination band system between 1.75 and 2.00\,$\mu$m and on the blue
edge of the 2.0\,$\mu$m H$_2$O band, respectively. Atmospheric
transmission at Cerro Paranal in these regions drops below 70--80\%
even at moderate precipitable water vapour
($\mathrm{PWV} \simeq 2.5$\,mm), and telluric correction with
\texttt{molecfit} \citep{Smette2015,Kausch2015} leaves residuals at the
few-percent level. On the broad ($\Delta v \sim 10^3$\,\kms) WC
emission profiles such residuals deform the line shape and bias the
CCF centroid by velocities comparable to or larger than the binary
detection threshold $\drvthresh$\,\kms. Spectroscopic monitoring of
WC9 stars by \citet{Williams2000} corroborates the unreliability of
NIR CCF centroids in this regime, and we therefore exclude the two NIR
lines from the primary classification.

\paragraph{\ion{He}{ii} variability and blending.}
The \ion{He}{ii}\,$\lambda 4686$ line forms higher in the accelerating
WC wind than the deeper \primaryline doublet
\citep{Hillier1989,Crowther2007}, making its profile more sensitive to
wind clumping, co-rotating interaction regions, and excess emission
from wind--wind collision in binaries. Stochastic profile variability
at the few-percent level on hour-to-day timescales is well documented
\citep{LepineMoffat1999}, and \citet{Shenar2019} explicitly recommend
avoiding \ion{He}{ii}\,$\lambda 4686$ for radial-velocity measurements
when alternatives exist. The \ion{He}{ii}\,$\lambda 5412$ line is in
addition blended with the \ion{C}{iv}\,$\lambda 5471$ feature in WC
spectra (Table~\ref{tab:emission_lines}), so that wind variability and
line confusion both contribute to a less stable centroid.

\paragraph{The \primaryline doublet -- primary RV line.}
The \primaryline doublet is the strongest and deepest-formed optical
emission feature in the WC spectrum and yields the most stable line
centroid in WC atmosphere modelling. We therefore adopt \primaryline
as the primary RV diagnostic for the binary classification of
Sect.~\ref{sec:criteria}. This choice also provides direct
methodological continuity with the precursor LMC WC binary survey of
\citet{Bartzakos2001}, which used the same line, and with the Galactic
WC survey of \citet{Dsilva2020,Dsilva2023}.

The remaining lines in Table~\ref{tab:emission_lines} -- the
\ion{O}{v}, \ion{O}{iv}, \ion{O}{vi}, \ion{C}{iv}\,$\lambda 3650$,
\ion{C}{iii}\,$\lambda 6700$ and \ion{C}{iv}\,$\lambda 7063$ features
-- are retained as cross-checks of the per-star RV time series
(Sect.~\ref{sec:line_agreement}), but are not used to flag binaries.

\subsection{Signal-to-noise requirements}
\label{sec:snr}

To ensure reliable radial-velocity measurements, we established minimum
signal-to-noise requirements per emission line through simulations. We
generated 30 synthetic single-lined binary (SB1) spectra using a WC4
template with varying S/N levels, measured the radial-velocity scatter for
each line, and required $\sigma_\mathrm{RV} < 3$\,\kms as the acceptance
threshold.

\subsection{Line-agreement cross-check}
\label{sec:line_agreement}

The choice of \primaryline as the primary RV line is supported by an
independent, data-driven cross-check across all \NlinesAll lines.

For each pair of emission lines $(\ell, \ell')$ we computed the Pearson
correlation coefficient $r_{\ell\ell'}$ of the per-star peak-to-peak
$\Delta\mathrm{RV}_{\max}$ values, restricted to the stars in which both
lines yielded valid measurements. To weight the correlation by sample
size, we define the agreement score per line $\ell$ as
\begin{equation}
S_\ell \;=\; \frac{\sum_{\ell' \neq \ell}
                    w_{\ell\ell'}\,r_{\ell\ell'}}
                  {\sum_{\ell' \neq \ell} 1},
\qquad
w_{\ell\ell'} \;=\; \frac{n_{\ell\ell'}}{\Nstars},
\label{eq:agreement_score}
\end{equation}
where $n_{\ell\ell'}$ is the number of stars with valid
$\Delta\mathrm{RV}_{\max}$ from both lines; we require
$n_{\ell\ell'} \geq \Nstarsmincorr$ for a pair to enter the average.
The line that is most ``agreeable'' with the rest of the sample has
the highest $S_\ell$.

The \primaryline doublet emerges as the highest-scoring line. As an
independent check, we computed the per-line equivalent width on the
same dataset and found a positive correlation between $S_\ell$ and the
mean per-epoch line strength: the strongest, most isolated WC emission
feature is therefore also the most kinematically self-consistent
across the sample. This is an empirical confirmation of the
depth-of-formation argument set out in Sect.~\ref{sec:primary_line}.
The agreement-score ranking and the score-vs-strength relation are
shown in Sect.~\ref{sec:results} (Fig.~\ref{fig:agreement}).

\subsection{$\Delta\mathrm{RV}$ threshold derivation}
\label{sec:threshold_derivation}

The binary detection threshold $\drvthresh$\,\kms used in
Sect.~\ref{sec:criteria} was derived empirically from the cumulative
observed binary fraction as a function of threshold.

\paragraph{Parametric Gaussian model.}
We model the cumulative binary fraction above threshold $T$ as the
survival function of a two-component Gaussian mixture, in which the
peak-to-peak $\Delta\mathrm{RV}_{\max}$ distributions of the
single-star and unresolved-binary populations are both approximated as
zero-centred normals,
\begin{equation}
f_\mathrm{bin}(>T)
  \;=\; (1 - f_\mathrm{bin})\,\Phi(T/\sigma_s)
        \;+\; f_\mathrm{bin}\,\Phi(T/\sigma_b),
\label{eq:gauss_threshold}
\end{equation}
where $\Phi$ is the survival function (i.e.\ $1 - \mathrm{CDF}$) of a
standard normal, $\sigma_s$ is the single-star peak-to-peak RV
scatter, and $\sigma_b$ is the corresponding scatter of the binary
population. The free parameters $(\sigma_s, \sigma_b, f_\mathrm{bin})$
are fitted to the observed $f_\mathrm{bin}(>T)$ vs $T$ curve by
non-linear least squares. The fitted parameters are
$\sigma_s = \sigmaSingleFit$\,\kms,
$\sigma_b = \sigmaBinaryFit$\,\kms, and
$f_\mathrm{bin} = \fbinAnalytic$, and the derived binary-detection
threshold is $T = \drvthresh$\,\kms, the value adopted throughout this
work.

\paragraph{Threshold value.}
The adopted threshold is the inflection of the modelled
$f_\mathrm{bin}(>T)$ curve: below it, the cumulative detection rate is
dominated by single-star wind variability, while above it, it is
dominated by genuine binaries. We adopt $\drvthresh$\,\kms as the
canonical binary-detection threshold (Sect.~\ref{sec:criteria}),
subject to the caveat below.

\paragraph{Caveat: the Gaussian assumption is approximate.}
Equation~\eqref{eq:gauss_threshold} assumes that both the single-star
and the unresolved-binary $\Delta\mathrm{RV}_{\max}$ distributions are
Gaussian. WR wind variability is, however, intrinsically stochastic
and asymmetric: clumping events, co-rotating interaction regions, and
line-profile excursions all produce non-Gaussian wings that the
analytic model of Eq.~\eqref{eq:gauss_threshold} cannot capture. We
have constructed a diagnostic tool that compares the observed
single-star $\Delta\mathrm{RV}_{\max}$ distribution to the best-fit
Gaussian via residual and second-derivative tests, and the diagnostic
confirms that the Gaussian is a useful first approximation but is not
exact in the wings. The bias-correction analysis described in
Sect.~\ref{sec:bias} is consequently formulated as a non-parametric
multinomial likelihood comparison (Eq.~\ref{eq:multinomial}) that does
not assume Gaussianity, and the threshold-curve cross-check in
Sect.~\ref{sec:threshold_fit} uses a Physics-based simulation rather
than an analytic Gaussian. A detailed treatment of the non-Gaussian
wind-variability distribution is deferred to a companion paper.

\subsection{Binary classification criteria}
\label{sec:criteria}

A star is classified as a spectroscopic binary if \emph{both} of the
following criteria are satisfied \citep[following][]{Bartzakos2001,
Dsilva2023}:

\begin{enumerate}
\item The peak-to-peak radial-velocity difference between any pair of
  epochs exceeds $\Delta\mathrm{RV}_\mathrm{max} > \drvthresh$\,\kms.
\item The detection is statistically significant:
  $\Delta\mathrm{RV} - \sigthresh\sigma > 0$, where
  $\sigma = \sqrt{\sigma_1^2 + \sigma_2^2}$ is the combined formal
  uncertainty of the epoch pair.
\end{enumerate}

The classification is performed in two stages. First, the epoch pair with
the maximum radial-velocity separation is evaluated. If this pair does not
satisfy both criteria, all remaining epoch pairs are scanned. The binary
decision is based exclusively on the \primaryline emission line, the
strongest and most isolated WC spectral feature
(Sect.~\ref{sec:primary_line}).

\subsection{Threshold-curve fit}
\label{sec:threshold_fit}

As a complementary check to the grid-based likelihood analysis described in
Sect.~\ref{sec:bias}, the cumulative observed binary fraction is fitted as
a function of the detection threshold $T$ on $\Delta\mathrm{RV}_\mathrm{max}$.
We adopt the Physics-based simulation: orbital parameters (period,
eccentricity, mass ratio, inclination) are drawn from the same priors
as in Sect.~\ref{sec:bias}, and the radial-velocity curve of each
synthetic binary is computed via Kepler's third law (Eq.~\ref{eq:K1})
and sampled at the observed cadence of each target. The free parameters
are the intrinsic binary fraction $f_\mathrm{bin}$, the single-star
scatter $\sigma_\mathrm{intrinsic}$, and the per-epoch measurement
uncertainty $\sigma_\mathrm{meas}$. Single stars are assigned intrinsic scatter at scale
$\sigma_\mathrm{intrinsic}$ and per-epoch noise at scale
$\sigma_\mathrm{meas}$, and the same two-criterion detection of
Sect.~\ref{sec:criteria} is applied. The cumulative binary fraction
$N_\mathrm{bin}(>T)/\Ntotal$ is computed across a grid of thresholds $T$
and compared to the observed cumulative curve to extract the best-fit
parameters.

\subsection{Spectroscopic morphology investigation}
\label{sec:spec_morphology}

In addition to the RV-based classification of Sect.~\ref{sec:methods}, each
star is visually inspected through a diagnostic-line viewer that groups the
spectrum into curated zoom panels covering interstellar features, telluric
bands, hydrogen Balmer and \ion{He}{i}/\ion{He}{ii} OB-type photospheric
absorption, the WR \ion{He}{ii}\,$\lambda 4686$ emission, the
\primaryline emission complex, and the
\ion{O}{v}/\ion{O}{vi} lines that distinguish WC from WO subtypes. Stars
showing morphological signatures inconsistent with a single WC/WO spectrum
-- for example, weak photospheric absorption components attributable to an
OB companion, or anomalous emission-line profiles -- are flagged as
candidate binaries. If any additional binaries are flagged through this
morphological inspection, the binary count is updated and the Monte-Carlo
bias correction is re-run with the revised \Nbin and \Ntotal. This
inspection adds \NbinSpec spectroscopically flagged binaries to the
RV-based count, raising the total to \NbinFinal/\Ntotal and the observed
binary fraction to $\fbinFinal$.

\section{Results}
\label{sec:results}

\subsection{Binary detections}
\label{sec:detections}

Applying the classification criteria described in
Sect.~\ref{sec:criteria} to the \primaryline radial velocities, we detect
\Ndetected new spectroscopic binaries among the \Nstars apparently single
targets. Combined with the \Nknown previously known binaries from
\citet{Bartzakos2001}, the total number of confirmed spectroscopic
binaries in the LMC WC/WO sample is \Nbin out of \Ntotal, yielding an
observed binary fraction of $\fbinobs$.

Table~\ref{tab:classification} lists the binary classification for each
star, including the maximum $\Delta\mathrm{RV}$, the associated
significance, and the epoch pair that triggered the detection.

\subsection{Radial-velocity measurements}
\label{sec:rv_results}

Per-star radial-velocity measurements, representative RV curves, and
example spectra spanning a clear SB1 detection, a marginal case at the
detection threshold, and an apparently single star are presented in
Sect.~\ref{sec:detections} and tabulated in
Appendix~\ref{app:rv_tables}.

\subsection{Observed binary fraction}
\label{sec:obs_fbin}

The observed binary fraction of $\fbinobs$ represents a lower limit on
the intrinsic binary fraction, as unfavourable orbital configurations
(long periods, low inclinations, low radial-velocity semi-amplitudes)
prevent detection for a fraction of the binary population. We address
this detection bias through Monte-Carlo simulations in
Sect.~\ref{sec:bias}.

\subsection{Threshold-curve fit results}
\label{sec:threshold_results}

The threshold-curve fit described in Sect.~\ref{sec:threshold_fit} is
applied to the observed cumulative binary fraction as a function of the
$\Delta\mathrm{RV}_\mathrm{max}$ detection threshold. The best-fit
parameters yield an intrinsic binary fraction
$f_\mathrm{bin} = \fbinThresh^{+\fbinThreshErrUp}_{-\fbinThreshErrDn}$,
with fitted single-star scatter $\sigma_\mathrm{intrinsic} = \sigmaIntrinsic$\,\kms
and per-epoch measurement uncertainty $\sigma_\mathrm{meas} = \sigmaMeasFit$\,\kms.
The threshold-curve estimate is compared to the grid-based likelihood
result of Sect.~\ref{sec:bias} as an internal consistency check on the
recovered intrinsic binary fraction.

\section{Bias correction}
\label{sec:bias}

The observed binary fraction is a lower limit on the true binary fraction,
since spectroscopic surveys are biased against detecting systems with long
orbital periods, low inclinations, or low mass ratios. Following the
methodology of \citet{Dsilva2023}, we employ a Monte-Carlo simulation to
constrain the intrinsic binary fraction $f_\mathrm{bin}$ and the
orbital-period distribution exponent $\pi$.

\subsection{Simulation framework}
\label{sec:sim_framework}

For each point on a two-dimensional grid in $(f_\mathrm{bin}, \pi)$, we
generate \Nsim synthetic stellar systems as follows:

\begin{enumerate}
\item Each system is assigned as binary with probability $f_\mathrm{bin}$,
  or single otherwise.

\item For binary systems, orbital parameters are drawn from the following
  distributions:
  \begin{itemize}
  \item \textbf{Period:} $p(\log P) \propto (\log P)^{\pi}$, with
    $\log P \in [0.15, 5.0]$\,days.
  \item \textbf{Eccentricity:} uniform on $[0, e_\mathrm{max}]$ with
    $e_\mathrm{max} = 0.9$.
  \item \textbf{Mass ratio:} $q = M_2/M_1$ uniform on $[0.1, 2.0]$.
  \item \textbf{Inclination:} $p(i) \propto \sin i$, $i \in [0, \pi/2]$
    (isotropic orientations).
  \item \textbf{Primary mass:} $M_1 = 10$\,\Msun (fixed) or drawn
    uniformly from $[10, 20]$\,\Msun.
  \end{itemize}

\item The radial-velocity semi-amplitude $K_1$ is computed from Kepler's
  third law:
  \begin{equation}
  K_1 = \left(\frac{2\pi G}{P}\right)^{1/3}
        \frac{M_2 \sin i}{(M_1 + M_2)^{2/3}}
        \frac{1}{\sqrt{1 - e^2}}.
  \label{eq:K1}
  \end{equation}

\item Radial-velocity curves are generated by solving Kepler's equation
  numerically, sampling at the actual observation epochs (MJD timestamps)
  of each of the \Nstars target stars. Crucially, the simulation is
  \emph{cadence-aware}: each simulated set contains exactly one star per
  real target, sampled at that target's actual MJD epochs. This faithfully
  reproduces the observational sampling and per-star cadence limitations.
  $N_\mathrm{sets}$ such sets are generated (typically $N_\mathrm{sets} =
  500$--3000), and the simulated CDF is constructed from the median binary
  fraction in each $\Delta\mathrm{RV}$ bin across all realisations.

\item Observational noise is added to \emph{both} single and binary
  stars. For single stars, the noise is drawn from a distribution with
  characteristic scale $\sigma_\mathrm{single}$ (intrinsic wind
  variability). For binary stars, per-epoch measurement noise is drawn
  from a separate error model distribution, ensuring that face-on or
  long-period binaries with $\Delta\mathrm{RV} \approx 0$ still exhibit
  realistic scatter. This prevents the likelihood score from favouring
  $f_\mathrm{bin} = 1$ (where noise-free face-on binaries mimic singles
  at zero cost). Both error model distributions may be fitted to the
  observed per-epoch RV uncertainties.
\end{enumerate}

\subsection{Comparison to observations}
\label{sec:likelihood_test}

For each grid point, we compute the peak-to-peak
$\Delta\mathrm{RV}_\mathrm{max}$ of each simulated star. A simulated star
is classified as a detected binary only if it satisfies the same two
criteria applied to the real data (Sect.~\ref{sec:criteria}): (1)
$\Delta\mathrm{RV}_\mathrm{max} > T$ for a given threshold $T$, and (2)
$\Delta\mathrm{RV} - \sigthresh\,\sigma_\mathrm{p2p} > 0$, where
$\sigma_\mathrm{p2p} = \sqrt{\sigma_\mathrm{max}^2 +
\sigma_\mathrm{min}^2}$ is the combined measurement uncertainty at the
epochs of maximum and minimum RV. This significance criterion prevents
noise-dominated detections at low thresholds and ensures consistency
between the simulated and observed binary fraction curves. The three known
Bartzakos binaries are added to the observed count at all thresholds
(denominator $N = 28$).

To compare the simulated and observed $\Delta\mathrm{RV}$ distributions,
we adopt the multinomial likelihood scheme of \citet{Dsilva2023}
(Sect.~4.2). The $\Delta\mathrm{RV}$ axis is partitioned into coarse
intervals with default edges $\{0,\,\drvthresh,\,250,\,650,\,\infty\}$\,\kms,
where the lowest edge coincides with the binary detection threshold so that
the first bin isolates non-detections and the remaining bins resolve the
high-amplitude tail. Let $n_i$ denote the observed count in bin $i$ and
$p_i$ the corresponding fraction predicted by the Monte-Carlo simulation
at a given grid point; the likelihood of the observed counts under the
simulated model is
\begin{equation}
\mathcal{L}(\mathbf{n}\,|\,\mathbf{p}) =
\frac{N!}{\prod_i n_i!}\,\prod_i p_i^{\,n_i},
\label{eq:multinomial}
\end{equation}
with $N = \Ntotal$. The best-fit model is the grid point that maximises
$\ln\mathcal{L}$.

\begin{figure}
\centering
\includegraphics[width=\columnwidth]{plots/cdf_obs_vs_sim.pdf}
\caption{Cumulative $\Delta\mathrm{RV}_{\max}$ distribution: observations
  (black solid step) vs.\ Monte-Carlo simulation at the best-fit grid
  point (red dashed; the 16--84\% percentile band of the simulation is
  shaded). The vertical dashed lines mark the multinomial bin edges
  (Eq.~\ref{eq:multinomial}). \emph{Preliminary snapshot}: the
  bias-correction grid is still being tuned and the values shown
  correspond to a representative grid configuration with
  $N_\mathrm{sim} = 1000$ realisations.}
\label{fig:cdf_obs_vs_sim}
\end{figure}

\subsection{Grid search and results}
\label{sec:grid_results}

The grid spans $f_\mathrm{bin} \in [0.3, 1.0]$ with step 0.005--0.01, and
$\pi \in [-3, +1]$. The search is parallelised across
$N_\mathrm{cores} - 1$ CPU cores.

The best-fit parameters are $f_\mathrm{bin} =
{\fbincorr}^{+\fbinerrup}_{-\fbinerrdn}$ and $\pi =
{\pibestfit}^{+\pierrup}_{-\pierrdn}$, where the uncertainties represent
68\% highest-density intervals (HDI68) from the marginalised
log-likelihood.

\begin{figure}
\centering
\includegraphics[width=\columnwidth]{plots/fbin_pi_heatmap.pdf}
\caption{Marginalised log-likelihood over $(f_\mathrm{bin}, \pi)$ for the
  Dsilva power-law period model. The cross marks the joint argmax; the
  white contour encloses the 68\% highest-density interval (HDI68). The
  likelihood is marginalised over the nuisance parameters
  $\sigma_\mathrm{single}$ and $\log P_\mathrm{max}$. \emph{Preliminary
  snapshot}: bias-correction grid not yet converged.}
\label{fig:fbin_pi_heatmap}
\end{figure}

\begin{figure}
\centering
\includegraphics[width=\columnwidth]{plots/fbin_pi_marginals.pdf}
\caption{One-dimensional marginal posteriors of (\textit{left}) the
  binary fraction $f_\mathrm{bin}$ and (\textit{right}) the
  period-distribution exponent $\pi$, derived from the Dsilva grid of
  Fig.~\ref{fig:fbin_pi_heatmap}. The solid vertical lines mark the
  joint argmax; the dashed lines and shaded regions mark the 68\% HDI
  bounds. \emph{Preliminary snapshot}.}
\label{fig:fbin_pi_marginals}
\end{figure}

\begin{figure}
\centering
\includegraphics[width=\columnwidth]{plots/bin_sensitivity.pdf}
\caption{Bin-choice sensitivity of the inferred binary fraction. For each
  multinomial-bin scheme tested, the marker is the argmax $f_\mathrm{bin}$
  and the horizontal bar covers the 68\% HDI. The default scheme of
  \citet{Dsilva2023} ($\{0,\drvthresh,250,650,\infty\}$\,\kms) is
  highlighted in red; alternative schemes vary in number of bins, edge
  placement, and edge spacing. The envelope across schemes is reported
  as a methodological-robustness statement alongside the headline value.}
\label{fig:bin_sensitivity}
\end{figure}


\section{Discussion}
\label{sec:discussion}

\subsection{Comparison with Galactic WC stars}
\label{sec:comparison_mw}

\citet{Dsilva2023} measured the observed spectroscopic binary fraction of
Galactic WC/WO stars using VLT/X-shooter spectroscopy and the same
cross-correlation framework adopted here, providing the natural
methodological twin of the present LMC measurement. Both surveys derive
radial velocities from the \primaryline emission line, classify a star as
a binary when $\Delta\mathrm{RV} > \drvthresh$\,\kms is recorded at
$\sigthresh\sigma$ significance, and apply Monte-Carlo bias correction to
recover the intrinsic binary fraction (Sect.~\ref{sec:bias}). The
methodological alignment removes the principal systematic differences
between published WR binary fractions, so the residual offset can be
attributed to the underlying stellar populations rather than to detection
choices. The cadence-dependent bias differs between the two samples: the
LMC targets were observed across \Nepochs epochs per star, whereas the
Galactic survey adopted per-target cadences tailored to expected orbital
periods. The corrected fraction \fbincorr, rather than \fbinobs (or
\fbinFinal), therefore provides the apples-to-apples comparison: a value
of \fbincorr above the Galactic estimate of \citet{Dsilva2023} would
indicate a stronger binary contribution at LMC metallicity, while a
comparable or lower value would constrain that prediction.

\begin{figure}
\centering
% \includegraphics[width=\columnwidth]{plots/lmc_vs_mw.pdf}
\caption{Comparison of the bias-corrected binary fraction and
  orbital-period distribution exponent for the LMC WC/WO sample (this
  work) and the Galactic WC/WO sample of \citet{Dsilva2023}. Error bars
  are 68\% HDIs from the marginalised log-likelihood. The horizontal
  dashed lines mark the Galactic reference values.}
\label{fig:lmc_vs_mw}
\end{figure}

\subsection{Metallicity dependence}
\label{sec:metallicity}

At the sub-solar metallicity of the LMC ($Z \approx 0.3$--$0.5\,\Zsun$;
\citealt{Russell1992}), line-driven mass loss is weaker than at solar
metallicity, raising the minimum initial mass at which a single star can
shed its hydrogen envelope and reach the WR stage \citep{Shenar2020}. A
generic prediction follows: as the single-star channel becomes
progressively inefficient, the binary channel must contribute a larger
fraction of the WC/WO population at LMC metallicity than at Milky Way
metallicity. Operationally, this expectation should manifest as
$f_\mathrm{bin,LMC} > f_\mathrm{bin,MW}$ once both samples are corrected
for observational bias. The relevant test is therefore a comparison of
the corrected LMC fraction \fbincorr with the corrected Galactic fraction
of \citet{Dsilva2023}, both of which marginalise over the assumed
orbital-parameter distributions. The corrected \fbincorr provides the
test: a value above the corresponding Galactic estimate would support
the metallicity-dependent binary-channel prediction, while a value
below or comparable to the Galactic estimate would constrain it.
A second observable carries complementary information about formation
channels: the orbital-period exponent \pibestfit recovered from the
bias-correction grid (Sect.~\ref{sec:bias}). A bottom-heavy period
distribution (more long-period systems) is expected when wide-orbit mass
transfer dominates the production of stripped helium stars, whereas a
top-heavy distribution (more short-period systems) points to
Roche-lobe-overflow stripping in close binaries. The two observables
together --- $f_\mathrm{bin}$ and the period-distribution shape --- are
therefore the primary handles on the metallicity dependence accessible
from this work.

\subsection{Formation channels}
\label{sec:formation}

Two competing channels are typically invoked to produce WC/WO stars: (a)
single-star evolution in which strong line-driven winds strip the
hydrogen and helium envelopes of an initially very massive star
\citep[the Conti scenario;][]{Conti1976}, and (b) binary mass transfer in
which Roche-lobe overflow strips the envelope onto a companion, exposing
the carbon-rich core at lower initial masses than the single-star
pathway requires \citep{Paczynski1967}.
Population-synthesis calculations predict that the binary stripping
channel becomes increasingly dominant as metallicity decreases, because
the single-star channel is throttled by the metallicity dependence of
line-driven mass loss \citep{Marchant2023}. The corrected binary fraction
\fbincorr is the principal observable that discriminates between the two
channels. A value $\gtrsim 50\%$ favours a binary-dominated origin for
the LMC WC/WO population, a value $\lesssim 30\%$ favours single-star
evolution, and intermediate values imply that both channels contribute
non-negligibly. The orbital-period exponent \pibestfit, together with
its Langer-model alternative (Sect.~\ref{sec:langer_model}), provides an
orthogonal test by tracing the imprint of the dominant interaction mode
on the surviving orbital architecture.
The mass-ratio distribution, which is in principle a third diagnostic,
is not constrained by this work: the survey is RV-only with limited
orbital-phase coverage per star (Sect.~\ref{sec:threshold_results}), so
the discussion here stops at the binary fraction and the
period-distribution shape.

\subsection{WC + compact-object candidates}
\label{sec:compact_objects}

As a follow-up to the RV-based binary classification, \NcandidateBH
stars in the sample have been flagged as candidate WC + compact-object
(WC+BH) systems and are being targeted for phase-resolved spectroscopic
monitoring (proposal in preparation; results will be appended to a
future companion paper). The initial selection was guided by the
orbital phase relationship between absorption and emission lines: in a
genuine WR+OB binary, photospheric absorption from the OB companion is
expected to phase-lock with opposing radial-velocity signatures
relative to the WR wind emission lines, while WR + compact-object
systems show either no or anomalous absorption-line phase behaviour.
The final candidate selection further refined this signature using a
complementary diagnostic that will be described in the companion paper.
WC+BH systems at LMC metallicity directly probe the black-hole mass
distribution from stripped helium-star progenitors, complementing
population-synthesis predictions of compact-object merger rates
\citep{vandenHeuvel2017}.

\section{Conclusions}
\label{sec:conclusions}

We have presented a multi-epoch VLT/X-Shooter spectroscopic multiplicity
survey of \Nstars apparently single carbon-rich Wolf--Rayet stars in the
Large Magellanic Cloud, combined with the \Nknown previously known
WC/WO binaries from \citet{Bartzakos2001}, for a total sample of
\Ntotal stars.

The radial-velocity analysis is anchored on the \primaryline emission
line --- the strongest and least variable WC wind feature, whose primacy
is confirmed independently by a line-agreement cross-check
(Sect.~\ref{sec:line_agreement}) showing that per-line RV
self-consistency correlates positively with emission-line strength.
Applying the joint $\Delta\mathrm{RV}_{\max} > \drvthresh$\,\kms and
$\sigthresh\sigma$ criterion of Sect.~\ref{sec:criteria}, we detect
\Ndetected new spectroscopic binaries, raising the total to \Nbin out of
\Ntotal and yielding an observed binary fraction of \fbinobs.
Following the cadence-aware Monte-Carlo bias-correction methodology of
\citet{Dsilva2023} and a multinomial-likelihood comparison on
$\Delta\mathrm{RV}_{\max}$ bins (Eq.~\ref{eq:multinomial}), the
intrinsic binary fraction is constrained to
$f_\mathrm{bin} = {\fbincorr}^{+\fbinerrup}_{-\fbinerrdn}$ and the
orbital-period-distribution exponent to
$\pi = {\pibestfit}^{+\pierrup}_{-\pierrdn}$ (Sect.~\ref{sec:bias}),
with 68\% highest-density intervals derived from the marginalised
log-likelihood.

Three orthogonal cross-checks bracket this result. The
two-component binary-evolution period model of \citet{Langer2020}
(Sect.~\ref{sec:langer_model}), the non-parametric threshold-curve fit
using a simulation-derived binary $\Delta\mathrm{RV}_{\max}$
distribution (Sect.~\ref{sec:threshold_fit}), and the bin-choice
sensitivity envelope across multiple $\Delta\mathrm{RV}$-binning
schemes are all tested for consistency with the headline values within
the 68\% HDI. To our knowledge, the bin-sensitivity analysis is the
first quantitative validation of binning robustness in a WR
multiplicity survey.

Several methodological caveats temper the headline values. The small
sample size ($N = \Ntotal$) places the multinomial likelihood in the
Poisson-noise-limited regime, so reported uncertainties are 68\% HDIs
from the joint posterior rather than Gaussian standard errors, and the
$\sim$3-semester observing baseline limits sensitivity to orbital
periods longer than $\sim 10$\,yr, leaving wide-orbit binaries as a
known incompleteness gap. The single-star RV scatter is poorly
constrained by the binned likelihood and is held fixed at a
representative value; the parametric Gaussian threshold derivation of
Sect.~\ref{sec:threshold_derivation} captures the bulk
$\Delta\mathrm{RV}_{\max}$ distribution but is approximate in the
wings, motivating the simulation-derived refinement of
Sect.~\ref{sec:threshold_fit}.

The bias-corrected fraction \fbincorr provides a direct empirical test
of the metallicity-dependent binary-channel hypothesis. A value above
the corresponding Galactic estimate of \citet{Dsilva2023} would support
the prediction that line-driven mass loss is too weak at LMC
metallicity for the single-star channel to dominate
\citep{Shenar2020, Marchant2023}, while a value comparable to or below
the Galactic value would constrain it. A definitive verdict awaits
finalisation of the bias-correction grid (Sect.~\ref{sec:bias}).

Three planned extensions of the present analysis will be reported in
companion papers: (i) a refined threshold-curve fit using a
non-Gaussian, simulation-derived parametric form for the binary
$\Delta\mathrm{RV}_{\max}$ distribution
(Sect.~\ref{sec:threshold_fit}), enabled by the converged
bias-correction parameters; (ii) photospheric-absorption confirmation
of all \Ndetected newly classified binaries together with a search for
hidden binaries among the apparently single stars
(Sect.~\ref{sec:spec_morphology}); and (iii) phase-resolved
spectroscopic monitoring of the \NcandidateBH candidate WC + compact-object
systems (Sect.~\ref{sec:compact_objects}) to constrain the black-hole
mass distribution from stripped-helium-star progenitors.


% ============================================================
% Acknowledgements
% ============================================================
\begin{acknowledgements}
Based on observations collected at the European Organisation for
Astronomical Research in the Southern Hemisphere under ESO programmes
\esoProgrammes (PI: T.\,Shenar).

G.S.\ acknowledges support from Tel Aviv University.
T.S.\ acknowledges support from [funding source].

This research has made use of the SIMBAD database, operated at CDS,
Strasbourg, France, and NASA's Astrophysics Data System Bibliographic
Services.
\end{acknowledgements}


% ============================================================
% References
% ============================================================
\bibliographystyle{aa}
\bibliography{references}

% ============================================================
% Appendix (online material)
% ============================================================
\begin{appendix}

\section{Radial-velocity measurements}
\label{app:rv_tables}

The full per-star, per-epoch radial-velocity measurements derived from the
\primaryline cross-correlation analysis of Sect.~\ref{sec:methods} are
available as supplementary material upon publication. For each of the
\Ntotal stars in Table~\ref{tab:sample}, the table reports the MJD of
observation, the measured radial velocity, the formal per-epoch
uncertainty $\sigma_\mathrm{RV}$ derived following \citet{Zucker2003},
and the peak-to-peak excursion
$\Delta\mathrm{RV}_\mathrm{max}$ used in the binary classification
(Sect.~\ref{sec:criteria}).

\section{Cross-correlation function profiles}
\label{app:ccf_plots}

\end{appendix}


\end{document}
